\section{Introduction}

How small can the total cross-sectional area of a finite cylinder covering
of a convex body be? Here a convex body $K\subset\R^3$ is compact, convex,
and has nonempty interior. A cylinder is a set
\[
 C=B+\R u,
\]
where $u$ is a unit vector and its base $B\subset u^\perp$ is measurable
with finite area. All areas and projections are Euclidean. Write
\[
 A_{\min}(K)=\min_{\dim E=2}|\pi_EK|,
\]
where $\pi_E$ is orthogonal projection onto the two-dimensional linear
subspace $E$ and $|\cdot|$ denotes planar Lebesgue measure.
The \emph{half-area cylinder covering conjecture} asks whether every
finite covering $K\subset\bigcup_{i=1}^m(B_i+\R u_i)$ satisfies
\begin{equation}\label{eq:conjectured-bound}
 \sum_{i=1}^m|B_i|\ge\tfrac12 A_{\min}(K).
\end{equation}

The cylinder question belongs to the line of covering problems that begins
with Tarski's plank problem. A \emph{plank} is the region between two
parallel hyperplanes, and its width is their distance. Bang proved that
a finite plank covering of a convex body has total width at least the
body's minimum width~\cite{Bang1951}; see also
\cite[Theorem~2.1]{Bezdek2009}.
Ball proved the directionwise refinement for centrally symmetric bodies:
after dividing each plank's width by the body's width in the same normal
direction, the sum is at least one~\cite[arXiv v1, Theorem~1]{Ball1991}.
Thus the width cost of a plank cover has both an absolute formulation
and a relative formulation adapted to the individual directions.

For cylinders in three dimensions, the two-cylinder cover of a regular
tetrahedron shows that replacing widths by base areas cannot preserve
the constant one. Its two axes are parallel to opposite edges, and its
total base area equals half the tetrahedron's minimum projection area.
Bezdek~\cite[Problem~3.1]{Bezdek2009} and
Bezdek and Litvak~\cite[arXiv v1, Introduction]{BL2009} attribute this example and
the resulting half-area question to Bang's 1951 paper. Their formulations
allow measurable bases perpendicular to the axes, as in our definition.

For a finite covering $\mathcal C=(B_i+\R u_i)_{i=1}^m$, define the
\emph{directionwise relative area} by
\begin{equation}\label{eq:relative-cost}
 \mathcal R(\mathcal C;K):=
 \sum_{i=1}^m\frac{|B_i|}{|\pi_{u_i^\perp}K|}.
\end{equation}
All denominators are positive because $K$ has nonempty interior.
Bezdek and Litvak proved $\mathcal R(\mathcal C;K)\ge1/3$
for every finite covering in three dimensions, and the stronger bound
$\mathcal R(\mathcal C;K)\ge1$ when $K$ is an
ellipsoid. In their proof, the general case combines the Rogers--Shephard
section--projection inequality with a volume comparison. For ellipsoids,
they reduce to a ball and use a density whose integral along every
interior chord is the same~\cite[arXiv v1, Theorem~3.1 and its proof]{BL2009}.
These arguments leave
a gap between the universal one-third lower bound and the proposed
one-half threshold.

Bezdek and Khan record the directionwise assertion
$\mathcal R(\mathcal C;K)\ge1/2$ under the name
\emph{1-Codimensional Cylinder Covering Conjecture}
\cite[arXiv v2, Definition~3 and Conjecture~4.13]{BK2016}.
It would imply \eqref{eq:conjectured-bound}, since
\begin{equation}\label{eq:cost-comparison}
 \mathcal R(\mathcal C;K)
 \le \frac{\sum_i|B_i|}{A_{\min}(K)}.
\end{equation}
Verreault's 2026 survey records the half-area question as
unresolved~\cite[Question~4.14]{Verreault2026}.
Our construction starts with the tetrahedral equality example itself
and replaces its two axes by finitely many nearby directions. It
violates both half-area bounds while retaining compact triangular bases.


We construct a counterexample with compact triangular bases. Set
$H=\sqrt2$ and
\begin{equation}\label{eq:tetrahedron}
 K=\{(x,y,Ht):0\le t\le1,\ |x|\le1-t,\ |y|\le t\}.
\end{equation}
Its vertices are $(\pm1,0,0)$ and $(0,\pm1,H)$, so it is a regular
tetrahedron of edge length $2$.

\begin{theorem}\label{thm:counterexample}
The tetrahedron \eqref{eq:tetrahedron} has $A_{\min}(K)=\sqrt2$.
For every $0<\eps\le1/2000$, it admits a covering by
$m=2\lceil2/\eps^2\rceil$ cylinders with compact triangular perpendicular
bases such that
\begin{equation}\label{eq:main-asymptotic}
 \frac1{\sqrt2}\sum_{i=1}^m|B_i|
 =\frac12-\frac{13}{6000}\eps^2+O(\eps^4).
\end{equation}
The remainder has absolute value at most $2\eps^4$, and the total base
area is strictly less than $A_{\min}(K)/2$.
\end{theorem}

Every covering in the theorem is finite, although its number of cylinders
increases as the tilt decreases. Similarities preserve the ratio of base
area to minimum projection area, so the construction applies to every
regular tetrahedron.

\begin{corollary}\label{cor:normalized-covering}
Every nondegenerate tetrahedron $T\subset\R^3$ admits a finite cylinder
covering with compact triangular perpendicular bases for which
\[
 \sum_{i=1}^m\frac{|B_i|}{|\pi_{u_i^\perp}T|}<\frac12.
\]
Thus the 1-Codimensional Cylinder Covering Conjecture is false.
\end{corollary}
\begin{proof}
For the regular tetrahedron $K$, each denominator is at least
$A_{\min}(K)>0$, so the sum is at most
$\sum_i|B_i|/A_{\min}(K)<1/2$ by
Theorem~\ref{thm:counterexample} with $\eps=1/2000$.

The directionwise ratios are invariant under invertible affine
maps~\cite[arXiv v1, Section~3]{BL2009}.
To see this directly, let $A$ be an invertible linear map and set
$v=Au/|Au|$. The linear map
\[
 L=\left.\pi_{v^\perp}A\right|_{u^\perp}:u^\perp\longrightarrow v^\perp
\]
is an isomorphism: $Lb=0$ implies that $Ab$ is parallel to $Au$,
hence $b$ is parallel to $u$, and therefore $b=0$.
The transformed cylinder has perpendicular base $LB$, and
$\pi_{v^\perp}(AK)=L(\pi_{u^\perp}K)$. Both areas are multiplied by
the same positive factor $|\det L|$, so their ratio is unchanged.
Translations also preserve it. Every nondegenerate tetrahedron is an
affine image of $K$; the images of its covering retain the same finite
count and compact nondegenerate triangular bases.
\end{proof}

This affine argument concerns the directionwise ratios. The conclusion
of Theorem~\ref{thm:counterexample}, which uses one minimum projection
area for all directions, is asserted here for regular tetrahedra.

\paragraph{Idea of the proof.}
Bang's equality example covers the parts $t\le1/2$ and $t\ge1/2$ by
cylinders parallel to the $x$- and $y$-axes, respectively. These axes
are parallel to opposite edges of the tetrahedron, and each base is a
triangle. We
subdivide these triangles into narrow sectors and give each sector its
own tilted axis. Tilting reduces the perpendicular area, but it can also
open gaps between neighboring sectors and between the two families.

Two choices prevent these gaps. First, the axes adjacent to a shared
sector side make the planes containing that side coincide. At a point
$(x,y,Ht)$ of the tetrahedron, these planes specify a finite list of
$y$-coordinates beginning at $-t$ and ending at $t$. Hence some adjacent
pair brackets $y$, even if the list is not ordered. This selects a sector
in the first family; exchanging the roles of $x$ and $y$ does the same
for the second. Second, we increase each sector's radial extent by a
quantity of order $\eps^2$. A gap between the families could occur only
near $t=1/2$. There the radial heights of the two intercepts, measured
from their respective tips, have opposite first-order shifts, $-\eps xy$ and $+\eps xy$, so their sum changes only
to second order. The increases can therefore be chosen small enough
that their area cost is less than the quadratic gain from tilting.

Section~\ref{sec:construction} proves the projection minimum and derives
the finite construction. Section~\ref{sec:coverage} proves exact coverage:
a first-crossing argument chooses angular sectors, and a single radial
budget excludes simultaneous failure of their cutoffs.
Section~\ref{sec:area} computes the perpendicular areas. The angular mesh
has size at most $\eps^2$, which keeps the total approximation error of
order $\eps^4$ even as the number of cylinders grows.

The same mechanism permits more freedom than the first construction
uses. Section~\ref{ext:cap-section} gives a common criterion for
continuous radial profiles, squared-radius cutoffs and bounded Cartesian
apertures, and compares midpoint sampling with endpoint-centred sectors.
Its estimates remain uniform as the partition is refined.
Section~\ref{sec:angular-cubic} then lets the radial margin tend to zero
with the tilt. In the same tetrahedron $K$, its normalized cost is
$1/2-\tau^2/240+(25/2)\tau^3+O(\tau^4)$: one family of constructions
attains the quadratic coefficient approached as the fixed margin tends
to zero. Its cylinders cover for $0<\tau\le1$, whereas the explicit
area bound guarantees strict saving for $0<\tau\le1/4000$.
Appendix~\ref{ext:alternate-estimates} gives complementary direct
proofs for curved caps, and Appendix~\ref{ext:scaling-specializations}
collects the exact counts and formulas under other parameter and scale
choices.
